\section{Conclusiones}
El presente trabajo cumplió con éxito el objetivo de desarrollar un sistema web enfocado en centralizar la carga, conversión, almacenamiento y gestión de documentos, imágenes y enlaces URL. Gracias a la integración con el mecanismo de autenticación de los sistemas de escritorio existentes en la institución, se garantiza el acceso exclusivo del personal autorizado y la seguridad de la información.
Con esta implementación, el PJETAM logra reducir los tiempos de búsqueda de documentos, eliminar el uso de herramientas externas y mejorar la accesibilidad a la información.

A nivel funcional, el sistema satisface por completo los requerimientos definidos para sus dos módulos. Por un lado, el módulo de conversión procesa adecuadamente documentos en distintos formatos para su conversión a PDF, y permite el almacenamiento de manera segura de imágenes y URLs. 
Por otro lado, el módulo de gestión documental centraliza la consulta, visualización, descarga y eliminación de archivos, en este destaca la implementación de códigos QR para cada archivo, funcionando perfectamente para la referenciación.

En la parte de desarrollo, el proyecto utilizó la arquitectura MVC, lo que permitió su elaboración de forma incremental y ordenada, primero se implementó la lógica de conversión de documentos a PDF, después se desarrolló el módulo de generación de códigos QR, y finalmente la integración con los servicios que necesitaban credenciales oficiales.

Como sugerencias para mejorar el sistema en un futuro, se identifican las siguientes: la implementación de limpieza automática de la carpeta local para no acumular archivos temporales, agregar paginación en la lista de gestión documental para mejorar la navegación cuando sean muchos archivos. Por último, sería conveniente integrar un algoritmo de compresión de imágenes antes de subirlas a S3 para optimizar el almacenamiento.

Desde una perspectiva de desarrollo personal, la elaboración de este proyecto fue de gran ayuda para aprender a integrar nuevas tecnologías y formas de llevar a cabo distintos procesos. Me permitió adquirir conocimientos en el uso de la arquitectura MVC. También apliqué herramientas nuevas para mí, como lo es Entity Framework para la gestión de la base de datos,
y puse en práctica conocimientos teóricos sobre AWS adquiridos en la universidad, comprobando sus ventajas y funciones en un escenario práctico. Para último, el trabajar en una institución real me brindó aprendizaje sobre la dinámica laboral, pude observar que el entorno laboral exige mayor autonomía y responsabilidad individual que en la escuela.
